\section{Einleitung}\label{einleitung}

Cloud-Dienste, Datenbanken oder das Internet – all diese Technologien basieren auf der Kommunikation zwischen verschiedenen Systemen. Die Interaktion erfolgt durch eine Client-Host-Architektur, die eine schnelle, stabile und effiziente Datenübertragung sowie die reibungslose Verarbeitung von Anfragen ermöglicht. In der modernen IT-Welt entfernen sich Systeme zunehmend von klassischen Serverräumen und entwickeln sich hin zu vernetzten Server- und Netzwerkarchitekturen. 

Die Bedeutung von verteilten und vernetzten Systemen nimmt stetig zu, da sie auf Netzwerken basieren, die die Grundlage der heutigen IT-Infrastruktur bilden. Der Begriff verteilte Systeme bezieht sich in diesem Zusammenhang auf eine Sammlung unabhängiger Einheiten, die miteinander kommunizieren, um ein größeres und effizienteres Gesamtsystem zu bilden. Dies lässt sich mit einer Firma mit verschiedenen Abteilungen vergleichen: Jede Abteilung arbeitet eigenständig, aber durch den kontinuierlichen Austausch von Informationen entsteht eine funktionierende Gesamtstruktur. 

Diese Konzepte bilden auch die Grundlage des Moduls "Vernetzte Systeme" (VNS), in dem wir die Möglichkeit hatten, uns intensiv mit der Gestaltung, Implementierung und Verwaltung verteilter Systeme auseinanderzusetzen. Ein zentraler Bestandteil dieser Veranstaltung ist der Einsatz von Docker-Umgebungen, die es ermöglichen, komplexe Anwendungen effizient in Containern zu organisieren und zu betreiben. 

Da Teamarbeit eine essenzielle Rolle in der modernen IT-Welt spielt, haben wir uns im Rahmen des VNS-Projekts zu einem Team zusammengeschlossen, um die erlernten Konzepte und Techniken in einem praxisnahen Projekt anzuwenden. Durch diesen praktischen Ansatz möchten wir nicht nur innovative Lösungen entwickeln und die Herausforderungen verteilter Systeme bewältigen, sondern auch unsere Problemlösungsfähigkeiten und Zusammenarbeit im Team weiter verbessern.

\section{Projektorganisation und Teamarbeit }\label{einleitung}
Für den Erfolg eines Projekts sind eine gute Organisation und eine effektive Zusammenarbeit im Team unerlässlich. Wir haben viele verschiedene Ideen gehabt, worüber unser Projekt handeln soll. Wir entschieden uns für eine simple ToDo-Liste, da wir uns hauptsächlich auf die verschiedenen Docker-Container fokussieren wollten, damit wir unser gelerntes über vernetzte Systeme und die Relevanz von Docker-Containern selbst und einer Virtual Machine auf die Probe zu stellen. Anschließend haben wir ein Git-Repository erstellt, um eine Basis für unsere Zusammenarbeit zu schaffen. Jedes Teammitglied hatte bestimmte Verantwortungen und Git ermöglichte es uns mehrere Baustellen parallel zu bearbeiten, um das Projekt effizient und zügig umzusetzen. Wir haben uns für eine flexible Arbeitsmethode entschieden und uns immer gegenseitig ausgeholfen, wodurch jeder an jedem Container und der Webanwendung mitgearbeitet hat. Wir haben aufeinander aufbauend gearbeitet und uns gegenseitig ergänzt, sodass jeder seine eigenen Ideen und Stärken einbringen konnte. In regelmäßigen Teamtreffen, welche sowohl Online als auch Präsenz stattfanden, konnten wir gemeinsam die Richtung des Projekts unter Konsens aller Mitglieder steuern und um eine kontinuierliche Kommunikation sowie einen reibungslosen Arbeitsablauf zu gewährleisten. Alle drei Tage haben wir uns über Jitsi getroffen, um den Fortschritt zu besprechen, Ideen auszutauschen und Probleme gemeinsam zu lösen. Diese regelmäßigen und kurzen Meetings halfen uns, fokussiert zu bleiben und Verzögerungen zu vermeiden. Zusätzlich dazu haben wir uns alle zwei Wochen in der Hochschule getroffen, um komplexere Themen zu diskutieren, den Code gemeinsam zu überprüfen und konkrete Pläne für die nächsten Schritte zu erstellen. Abseits der Teammeatings haben wir uns per Chat-Messenger stets über Änderungen und potenzielle Anregungen oder Ideen informiert, da uns eine offene Kommunikation und eine transparente Arbeitsweise enorm wichtig war. Durch diese klare Struktur und regelmäßigen Meetings konnten wir das Projekt effizient verwalten und sicherstellen, dass alle Teammitglieder aktiv beteiligt waren. Die Kombination aus flexiblen Online-Meetings und persönlichen Treffen ermöglichte uns eine perfekte Balance zwischen Produktivität und Zusammenarbeit. Außerdem haben wir unseren Fortschritt sukzessiv dokumentiert, was es uns erleichterte, den Überblick über erledigte und noch offene Aufgaben zu behalten. Am Ende half uns dieser strukturierte Ansatz, das Projekt erfolgreich und ohne größere Hindernisse abzuschließen.
